\section{Motivación, Datos e Hipótesis Estructural}

% ================================================================
\begin{frame}{El Dilema Regional ante el Choque COVID-19 (SCQA)}
\vspace{-0.18cm}
\begin{columns}[T,onlytextwidth]
\begin{column}{0.51\textwidth}
\begin{contrastbox}{ccdeblue}{Marco Estructural SCQA}
\scriptsize\RaggedRight
\begin{itemize}\setlength{\itemsep}{0.06em}
  \item \textbf{Situación ($S$):} Los 24 departamentos exhiben fuerte heterogeneidad productiva, desde enclaves mineros ($>40\%$) hasta economías urbanas.
  \item \textbf{Complicación ($C$):} En 2020 el PIB colapsó, pero los modelos en tasas anuales ($\Delta\%\text{PIB}$) sugerían un impacto indiferenciado ($p > 0.20$).
  \item \textbf{Pregunta ($Q$):} ¿Fue el choque neutral a la estructura productiva o las tasas anuales ocultan una cicatriz acumulada asimétrica?
  \item \textbf{Respuesta ($A$):} Al aplicar CoDa (CLR) y pasar a \textbf{niveles indexados ($\text{PIB}_{2019=100}$)}, emerge un rezago pospandemia significativo al $1\%$ en los enclaves.
\end{itemize}
\end{contrastbox}
\end{column}

\begin{column}{0.46\textwidth}
\begin{ccdebox}[left=2mm,right=2mm,top=1mm,bottom=1mm]
\scriptsize
\textbf{\textcolor{ccdenavy}{Arquitectura de la Base Consolidada (INEI)}}
\vspace{0.04cm}

\resizebox{\linewidth}{!}{%
\renewcommand{\arraystretch}{1.10}
\begin{tabular}{@{}l r@{}}
\toprule
\textbf{Dimensión del Panel} & \textbf{Magnitud} \\
\midrule
Unidades transversales ($N$) & \textbf{24 departamentos} \\
Horizonte temporal ($T$) & \textbf{2008--2025 (18 años)} \\
Observaciones totales & \textbf{428 filas de panel} \\
Ventana pre-shock (Clustering) & \textbf{2015--2019 (5 años)} \\
Ventana ciclo COVID (Niveles) & \textbf{2017--2023 (100\% bal.)} \\
Ramas de actividad VAB & \textbf{12 sectores (suma 100\%)} \\
\bottomrule
\end{tabular}%
}
\end{ccdebox}

\vspace{0.05cm}
\begin{keybox}[left=2mm,right=2mm,top=0.9mm,bottom=0.9mm]
\scriptsize\RaggedRight \textbf{So What?}: Clasificar con el promedio 2015--2019 garantiza \textbf{exogeneidad estricta} frente al shock de 2020.
\end{keybox}
\end{column}
\end{columns}
\source{INEI -- Cuentas Nacionales Departamentales (Cuadros 1 y 2, precios constantes de 2007).}
\end{frame}

